\section{Eurosystem purchases of sovereign bonds}\label{sec:eurosystem-purchases}

Guarantees support repayment after the sovereign shock. Bond purchases instead change banks' exposure before it occurs. We consider Eurosystem purchases of bank-held sovereign bonds at the model's date-1 book value of one, with selling banks receiving an equal amount of safe reserves. Such purchases create reserves on bank balance sheets \citep{ECBReserves2021}.

Each experiment specifies the sovereign issuers whose bonds are purchased and a common purchase fraction. For each selected issuer, the Eurosystem buys the stated fraction of the bonds held by each of the four country banking blocks, so purchases across banks are proportional to existing holdings. For example, purchasing 50\% of French holdings means buying half of the French sovereign bonds held by each German, French, Italian and Spanish bank block, while leaving their other sovereign holdings unchanged. In the all-issuer experiments, the stated fraction applies to holdings of all four sovereigns.

Let $P_{ig}\in[0,x_{i,G_g}]$ be the amount of sovereign $g$'s bonds purchased from bank $i$. The bank's direct resource gain in state $s$, holding other banks' repayments fixed, is
\[
 \Delta A_i(s)=\sum_gP_{ig}[1-R_g^G(s)].
\]
Replacing a defaulting sovereign claim with reserves reduces the bank's initial loss. The resulting change in interbank receipts can also remove adverse clearing outcomes. Sovereign debt outstanding is unchanged: the Eurosystem holds the sold bonds and issues the corresponding reserve liabilities. Bank total assets, debt, accounting equity, headroom and private-credit losses remain fixed; prices and credit supply do not respond. Appendix~\ref{app:purchase-implementation} gives the allocation and balance-sheet details.

\subsection{Purchase composition and banking outcomes}

Table~\ref{tab:eurosystem-purchases} compares purchase sizes and issuer compositions. For each experiment, we solve the altered network in every sovereign state and recompute the minimum repayment-floor capacity.

\begin{table}[htbp]
\centering
\floatcaption{Sovereign bond purchases and banking equilibria}{Amounts are EUR billions. Purchases replace bank-held sovereign bonds with safe reserves before the shock at book value. Defaults under full interbank repayment are evaluated after sovereign and private-credit losses, assuming that every interbank claim is paid in full; the two state-count columns range from zero to sixteen. Equilibria is the total across those states. Floor capacity is the attained minimum common capacity under proportional rationing for the post-purchase network; a dash denotes zero. Purchase principal is an asset exchange, not a fiscal cost. All other calibration inputs are unchanged.}
\label{tab:eurosystem-purchases}

\small
\renewcommand{\arraystretch}{1.12}
\begin{tabularx}{\textwidth}{@{}>{\raggedright\arraybackslash}Xrrrrr@{}}
\toprule
Purchased holdings & \shortstack{Purchase\\principal} & \shortstack{Default states\\under full\\interbank repayment} & \shortstack{Multiple\\states} & Equilibria & \shortstack{Floor\\capacity} \\
\midrule
No purchases & $-$ & 12 & 8 & 25 & 60.6 \\
All issuers: 25\% & 292 & 12 & 6 & 22 & 14.2 \\
All issuers: 50\% & 585 & 12 & 0 & 16 & $-$ \\
All issuers: 75\% & 877 & 0 & 0 & 16 & $-$ \\
France only: 50\% & 151 & 12 & 0 & 16 & $-$ \\
Italy and Spain: 75\% & 540 & 0 & 8 & 27 & 39.6 \\
\bottomrule
\end{tabularx}


\end{table}

Purchasing 25\% of banks' sample-country sovereign holdings, EUR~292 billion, reduces multiple states from eight to six and minimum floor capacity from EUR~60.6 billion to EUR~14.2 billion. Purchasing 50\%, EUR~585 billion, eliminates multiplicity while defaults under full interbank repayment remain in twelve states. At 75\%, EUR~877 billion, every banking equilibrium is fully performing.

Issuer composition gives the sharper distinction. Half of banks' French sovereign holdings, EUR~151 billion, suffices to eliminate multiplicity without changing the defaults under full interbank repayment. These purchases weaken the French repayment feedback supporting adverse outcomes. Purchasing three quarters of Italian and Spanish holdings, EUR~540 billion, instead makes full bank repayment an equilibrium in every sovereign state, but leaves eight states with additional equilibria in which banks default. In those adverse outcomes, reduced interbank receipts sustain the defaults, including through domestic interbank claims. Minimum floor capacity remains EUR~39.6 billion. Making full repayment possible and eliminating adverse equilibria are distinct objectives; equilibrium counts need not decline monotonically with purchases.

\subsection{Risk transfer and policy interpretation}

The bank resource gain is matched by a cash shortfall on the Eurosystem's new assets relative to its reserve liabilities. With half of French holdings purchased, the largest additional shortfall across sovereign states is about EUR~75 billion. This is a contingent public exposure, not an expected fiscal cost. The model does not allocate Eurosystem losses among national central banks or feed their financing back to households and sovereigns.

Purchases can therefore interrupt repayment feedback, but transfer sovereign risk to the central bank. The repayment floor can instead eliminate adverse outcomes without equilibrium transfers. Purchase principal and guarantee capacity measure different commitments and cannot establish a fiscal-cost ranking. The exercise isolates exposure reallocation, complementing the valuation and lending channels discussed in the introduction.

